% Generated draft: requires mathematical and source review.
\documentclass[11pt]{article}
\usepackage[margin=25mm]{geometry}
\usepackage{fontspec}
\setmainfont{Latin Modern Roman}
\usepackage{amsmath,amssymb,mathtools,unicode-math}
\setmathfont{Latin Modern Math}
\usepackage{enumitem,needspace,xurl,hyperref}
\setlength{\parindent}{0pt}
\setlength{\parskip}{5pt}
\setlength{\emergencystretch}{3em}
\newcommand{\N}{\mathbb{N}}
\newcommand{\Z}{\mathbb{Z}}
\newcommand{\Q}{\mathbb{Q}}
\newcommand{\R}{\mathbb{R}}
\newcommand{\C}{\mathbb{C}}
\newcommand{\F}{\mathbb{F}}
\newcommand{\A}{\mathbb{A}}
\newcommand{\PP}{\mathbb{P}}
\newcommand{\E}{\mathbb{E}}
\newcommand{\eps}{\varepsilon}
\newcommand{\dd}{\,\mathrm{d}}
\begin{document}
\section*{UnsolvedMath ID: 30000691}
\textbf{Source problem code:} catalogue entry
\subsection{Definitions and mathematical statement}
\textbf{Extremal neighborliness problem.} How neighborly can a centrally symmetric polytope be as a function of its dimension and number of vertices? For a full-dimensional \(d\)-polytope, translate its symmetry center to zero, so \(P=-P\). Its vertices occur in antipodal pairs \(v,-v\); hence the vertex count \(N\) is even, with the usual admissible range \(N\ge2d\).\par In the central-symmetry-adapted convention, \(P\) is \(k\)-neighborly if every subset of at most \(k\) vertices containing no antipodal pair is the vertex set of a proper face. A precise extremal target is\[\nu(d,N)=\max\{k\in\{0,\ldots,d\}:\exists\text{ centrally symmetric }d\text{-polytope with }N\text{ vertices having this property}\}.\] Determine this function, or sharp asymptotic bounds in stated regimes of \(d,N\), together with constructions achieving the proposed order. The explicit range \(k\le d\) avoids a vacuous unbounded parameter when only a limited number of nonantipodal vertices can be selected. No asymptotic candidate or required approximation factor is supplied.\par The source does not define neighborliness. Excluding antipodal pairs is essential to the intended centrally symmetric version: an antipodal segment passes through the interior and cannot be an edge, so ordinary two-neighborliness would already fail. The adapted convention above is therefore explicit and requires source confirmation, not silently confused with ordinary neighborliness of unrestricted polytopes. The source asks for dependence on both dimension and vertex count, not merely examples with a fixed number of vertices.\par The supplied triage checked 2026-08-21 labels the sharp asymptotic problem open, reporting substantial constructions and upper bounds without a sharp general answer. It supplies no exact quantitative theorem in the record. This dated assessment requires independent verification.
\subsection{Short English statement}
Determine the maximum central-symmetry-adapted neighborliness for each admissible dimension and vertex count, or its sharp asymptotic order. Nonantipodal vertex subsets are the relevant ones; ordinary neighborliness would be blocked immediately by antipodal pairs.
\subsection{Sources}
\begin{itemize}\item Local UnsolvedMath registry, permanent ID 30000691, code catalogue entry; \emph{Neighborliness of Centrally Symmetric Polytopes}, Oberwolfach Reports Open Problems.\item \emph{Geometric and Topological Combinatorics (2007)}, Oberwolfach Reports, DOI 10.4171/owr/2007/04, pp. 195--272. \url{https://doi.org/10.4171/owr/2007/04}. Proposed year 2007; no proposer named.\item Supplied literature label \emph{original problem report}: \url{https://doi.org/10.4171/owr/2007/04}.\item Supplied literature label \emph{Novik--Ziegler constructions}: \url{https://arxiv.org/abs/math/0507280}.\item Supplied literature label \emph{explicit constructions and bounds}: \url{https://arxiv.org/abs/1203.6867}.\item Literature triage checked 2026-08-21, status \texttt{open}; record updated 2026-09-16. Statement status \texttt{exact}, reported retained byte-for-byte. Verify dated claims independently.\end{itemize}
\end{document}
